\chapter{Rendering on slow links}
\label{ch:slowlinks}

\begin{aside}
Your app works locally. Then a user runs it over SSH from a
coffee-shop wifi and the screen lags by seconds. The diff was
small; the link is the problem. There are mitigations.
\end{aside}

\section{Where bandwidth goes}

A typical TUI sends:

\begin{itemize}[leftmargin=*]
  \item Initial alt-screen entry: ~100 bytes.
  \item Full-screen first paint: 1--10 KB depending on
    content and colour density.
  \item Steady-state frame: 30--500 bytes per active frame.
  \item Idle: zero.
\end{itemize}

Over a 1 Mbps SSH session that is one frame per ~5 ms; over
56 kbps dial-up (some still exist) it is 50 ms per KB.

\section{Reducing wire bytes}

\begin{itemize}[leftmargin=*]
  \item Avoid full-screen redraws. Keep diffs minimal.
  \item Cache styles aggressively (don't re-emit unchanged
    SGR).
  \item Use \code{ESC [ K} (erase to end of line) instead of
    spaces when a row's tail is uniform.
  \item Avoid OSC 8 hyperlinks on every cell; share IDs.
\end{itemize}

\maya{} does the first three automatically. The fourth is the
app's responsibility.

\section{Mosh}

For very slow or intermittent links, the Mosh project
(\code{mosh.org}) provides a Predict-and-Sync layer: client
predicts what the screen should look like and the server
sends only deltas to the client's local prediction.

\maya{} apps work over Mosh transparently. The wire savings
are dramatic on flaky links.

\section{Reducing colour}

Truecolour is 18--22 bytes per cell. 256-colour is 11. Plain
text is 0 (no SGR). On a slow link, a flag like
\code{--mono} that disables colour can be a usability win.

\section{Lazy rendering}

For huge logs, don't render the parts the user cannot see.
\maya{}'s scroll widget already does this: only the visible
viewport is painted. The data structure is the full log;
the painted slice is the small window.

\section{Latency vs throughput}

A bandwidth-constrained link delivers many bytes per second
eventually. A high-latency link makes every small frame feel
slow. The two need different fixes:

\begin{itemize}[leftmargin=*]
  \item Low bandwidth: shrink frames.
  \item High latency: predict locally (Mosh).
\end{itemize}

\section{Synchronised output}

DEC 2026 reduces visible mid-frame artifacts. On a slow
link, mid-frame visibility is amplified --- the user sees
half a frame for hundreds of milliseconds. \maya{} uses
2026 when available.

\section{Recap}

\begin{recap}
\begin{itemize}[leftmargin=*]
  \item Minimal diffs help everyone, especially slow links.
  \item Erase-to-end-of-line beats trailing spaces.
  \item Drop colour for severe constraint.
  \item Mosh hides latency at the protocol layer.
  \item Synchronised output reduces visible mid-frame
    artifacts.
\end{itemize}
\end{recap}

\section*{Workshop}

\begin{enumerate}[leftmargin=*]
  \item SSH into a remote box and run your app. Measure
    bytes per frame via \code{tcpdump}.
  \item Try Mosh; compare typing latency.
  \item Implement a \code{--mono} flag and measure bandwidth
    reduction.
\end{enumerate}
